%\documentclass[11pt]{article}
\documentclass[12pt,bibliography=numbered]{article}
\usepackage{pdfpages}
\usepackage{wrapfig}
\usepackage{longtable}
\usepackage{xcolor}
\usepackage{colortbl}
\usepackage{amsmath}
\usepackage{amsfonts}
\usepackage{paralist}
\usepackage[top=0.5in, bottom=1.0in, left=0.75in, right=0.8in]{geometry}
\usepackage{graphicx}
\usepackage[sort&compress,numbers]{natbib}
\usepackage{url}
\usepackage{color}
\usepackage{xspace}
\usepackage{caption}
\usepackage{subcaption}
\usepackage{adjustbox}
\usepackage{enumitem}
\usepackage{comment}
\usepackage{hyperref}[colorlinks=true,linkcolor=blue,citecolor=red]

\usepackage{endnotes}
\renewcommand{\footnote}{\endnote}

\renewcommand{\enotesize}{\footnotesize}
\renewcommand\enoteformat{%
  \raggedright
  \leftskip=1.0em
  \makebox[0pt][r]{\theenmark. \rule{0pt}{\dimexpr\ht\strutbox}}%
}
%\usepackage[hang, flushmargin, ragged]{footmisc}

\usepackage{fancyvrb}

%\usepackage[utf8x]{inputenc} 
%
%\DeclareUnicodeCharacter

\newcommand{\TODO}[1]{\textcolor{red}{\textbf{TODO} #1}}
\newcommand{\CHECK}[1]{\textbf{\textcolor{red}{CHECK: #1}}}
\newcommand{\DRAFT}[1]{\textbf{\textcolor{blue}{DRAFT:}}\textcolor{blue}{ #1}}


\renewcommand{\familydefault}{\sfdefault}

\begin{document}
\title{Una-May O'Reilly\\Professional Service
}
\author{  }

\date{\today}

\maketitle

\section*{Professional Service}

Assuming leadership in different communities fulfills my passion for service.  The scientific field of Evolutionary Computation is fairly youthful and its sub-field of Genetic Programming even younger. Since I was a PhD student, I have led different efforts to expand these fields and make them more accessible. I chaired the first workshop for graduate students at Genetic Programming 1997. When I was relatively junior, in 2005, I served as the first female chair of ACM SIGEVO's GECCO, the most prestigious conference of the field. I also was the first female co-chair of the European Conference on Genetic Programming. I was an inaugural member of the executive board of ACM SIGEVO and continue to serve as a member, having served at different times as an Officer, i.e. Secretary and Vice Chair. I am an inaugural member of the editorial board of the Springer journal Genetic Programming and Evolvable Machines, and currently am the journal's Area Editor for Data Analytics and Knowledge Discovery.  The ACM Journal on Transactions on Evolutionary Learning and Optimization was founded just recently and I serve as an Area Editor for it. 
I annually teach the Introduction to Genetic Programming tutorial at GECCO and recently have started to offer a tutorial on Coevolutionary Computation for Adversarial Deep Learning at GECCO.  In 2013 a career recognition award propelled me to shift prevailing culture by starting an annual meeting for Women in Evolutionary Computation. To this day, it is  regularly hosted at  GECCO.  I am especially proud to report that the meeting's scope has been updated to include any under-represented member of the EC community.   
%I deliberately try  to live my life locally -- being mindful of the present, those around me, and my role relative to those closer to me.
%\begin{comment}

In 2013 I updated my research group’s name to ALFA from Evolutionary Design and Optimization. Based on its application focus on security, ALFA is mostly funded by DARPA and the US Government, with additional funding from multiple CSAIL Alliance Program initiatives.  %ALFA stands for \textit{Any Scale Learning For All.}  The "Any Scale Learning" refers to wide ranging technical activities such as
%working on, not only statistical machine learning, but also machine learning inspired by biological evolution, self-styled \textit{evolutionary algorithms}.
%., and
%working on machine learning algorithms for not only large data but also data that is sparse in various dimensions.
%For data at large and small scales, we develop customizations and new machine learning algorithms. 
%We design systems not only for workstations, but also for the cloud and super computers. 
%We consider not only single agents, but also multiple agents, and not only single agent interactions, but also multi-agent interactions.
%The "For All" part conveys ALFA's commitment to research questions arising from problems that allow ALFA to contribute to socially relevant domains, such as cyber security.
%of ALFA's name 
Among the many rewards of leading ALFA, my group serves my passion for mentoring students and postdocs from MIT. ALFA has been a humane academic research home to numerous UROPs, Masters students (in Computer Science, Technology and Policy, and System Design \& Management), PhD students, and PostDocs. For MIT's Department of Electrical Engineering and Computer Science, I have led the Machine Learning track for PhD admissions, in addition to serving on the admissions committee multiple years.  I also enjoy academically advising EECS undergraduates and Ph.D. students.   

Besides working for technical and academic communities, I also devote efforts to the CSAIL community. I have led the CSAIL Pretty Committee, Lobby Committee, and Logo Committee. I lead the Applied Machine Learning Community of Research and I led the committee organizing CSAIL's Celebration of its  20th and 60th anniversaries.  In the service of principal and senior research scientists in CSAIL, I have worked on defining the appointments' review process. I am currently working on recommendations to CSAIL that will improve the transparency and details of the appointments' responsibilities and career paths.  A detailed, full list of my service can be found in my Curriculum Vitae.  

%\end{comment}

%  funding
%RIVALS is  largely funded by DARPA (XD3 and CHASE programs, CASTLE program in submission) and the US Government.
%
%BRON development continues and it is largely funded by DARPA (CHASE program  and CASTLE program in bidding) and the US Government.

%I have reviewed proposals for National Science Foundation (USA), Science Foundation of Ireland, Science Foundation of Israel, Science Foundation of Switzerland, Natural Sciences and Engineering Research Council, Canada, and Science Foundation of the Netherlands.

%In 2003, I filmed a segment on the DragonFlyTV show, “Real Scientists” which appears on PBS Kids Go!  



\end{document}